% =============================================================
%  SECTION 2 — RELATED WORK  (revised)
%  Reduction target: ~40% shorter than original
%  Structure: 4 subsections, each following the pattern:
%    what prior work did → remaining limitation → how this work differs
%  All textbook background on CNNs, transformers, sampling removed.
%  "Macro-context" macro-economics paragraphs deleted.
% =============================================================
\section{Related Work}

\subsection{Document Understanding and OCR-Based Information Extraction}

Layout-aware document understanding has advanced considerably with the
introduction of LayoutLM~\cite{xu2020layoutlm} and its successors, which
jointly encode text tokens and their two-dimensional bounding-box coordinates
using transformer self-attention. These models achieve strong performance on
form understanding and document classification benchmarks, but require
fine-tuning on annotated data in the target layout domain. End-to-end
approaches such as Donut~\cite{kim2022ocr} eliminate the OCR stage entirely
by learning to decode structured outputs directly from document images;
however, they are computationally expensive at inference time and require
substantial task-specific training data.

For enterprise RPA, where document formats vary across suppliers and input
quality fluctuates between high-resolution scans and low-quality phone
photographs, both categories of model face practical constraints. Our pipeline
delegates visual understanding to a cloud OCR service with layout analysis
and isolates the semantic extraction step in a downstream LLM, decoupling
visual-quality sensitivity from extraction logic without requiring fine-tuning.

\subsection{Vietnamese Automatic Speech Recognition}

Vietnamese presents specific difficulties for ASR: it is a tonal,
monosyllabic language with six tone classes, limited publicly available
training data, and pronounced regional dialect variation. Wav2Vec
2.0~\cite{baevski2020wav2vec} demonstrated that self-supervised pre-training
on unlabelled speech can significantly reduce the labelled data requirement
for low-resource languages, including Vietnamese. The Zipformer
architecture~\cite{yao2024zipformer} further improved parameter efficiency by
processing speech at multiple temporal resolutions.

PhoWhisper~\cite{le2024phowhisper}, developed by VinAI, fine-tunes the
OpenAI Whisper encoder-decoder~\cite{radford2023whisper} on over 800 hours
of multi-regional Vietnamese speech and achieves lower WER than
Wav2Vec-based systems on conversational Vietnamese. VietASR~\cite{zhuo2025vietasr}
combines a HuBERT~\cite{hsu2021hubert} encoder with a Zipformer decoder to
maintain competitive accuracy while supporting real-time streaming. A
limitation shared by both systems in the retail domain is that pre-training
corpora do not cover the mixture of Vietnamese, English product names, and
abbreviated trade terminology characteristic of spoken purchase orders.

This work selects PhoWhisper-base as the acoustic transcription module and
applies an LLM post-correction layer specifically targeted at known
phonetic substitution patterns in retail speech (e.g., the substitution of
homophonic loanwords for technical product names).

\subsection{Zero-Shot Information Extraction via LLMs}

Traditional named entity recognition (NER) systems trained on annotated
corpora—such as BiLSTM-CRF models~\cite{lample2016neural}—require
substantial labelled data per entity schema and do not generalise when new
entity types or document layouts are introduced. Prompt-based extraction
using LLMs bypasses this requirement: given a schema description and a few
illustrative examples, models such as GPT-4 and Gemini can extract structured
records from free-form text in a zero-shot or few-shot
setting~\cite{reid2024gemini}.

The practical constraint is throughput. Invoking a frontier LLM for each
document in a high-volume pipeline produces per-document latency in the
range of one to several seconds and incurs API costs that scale linearly with
token count. Our architecture limits LLM invocation to two narrow functions—
entity extraction from transcribed text and optional taxonomy prediction—
rather than using it for end-to-end document processing, thereby bounding
both latency and cost.

\subsection{Hybrid Retrieval for Entity Resolution}

Sparse retrieval with BM25~\cite{robertson2009probabilistic} provides fast,
exact-match recall but fails on synonyms and abbreviations. Dense retrieval
with sentence embeddings~\cite{reimers2019sentence} captures semantic
similarity but can confuse closely related product specifications that differ
only in a numeric attribute (e.g., \texttt{LED 8W} vs.\ \texttt{LED 9W}).
Hybrid approaches that combine both signals through Reciprocal Rank
Fusion~\cite{cormack2009reciprocal} have been shown to outperform either
strategy in isolation on heterogeneous retrieval benchmarks.

For Vietnamese product catalogs, which contain both semantically descriptive
product names and opaque alphanumeric SKU codes, neither modality is
universally preferable. We implement RRF as the primary fusion strategy within
the Qdrant vector database~\cite{qdrant}, which supports parallel execution of
dense and sparse retrieval in a single query. In addition, we evaluate two
variants: (i) \emph{Weighted Score Fusion} (WSF), which applies min-max
normalisation before a linear combination of scores with tunable weight
$\alpha$, and (ii) \emph{Weighted RRF} (WRRF), which applies $\alpha$
directly to the rank-based contributions, eliminating the normalisation
requirement. The three strategies are compared empirically in
Section~\ref{sec:experiments}.

A further distinction from prior retrieval work is the optional
taxonomy-guided filter applied before retrieval. By conditioning the vector
search on an LLM-predicted product category label, the effective catalog size
presented to the retrieval engine is reduced, lowering the risk of
cross-category semantic collisions in large-scale catalogs.